\documentclass[10pt,a4paper]{amsart}
\usepackage[margin=19mm]{geometry}
\usepackage{amsmath,amssymb,mathtools}
\usepackage[scr=rsfs,cal=euler]{mathalfa}
\usepackage{xfrac}
\usepackage[unicode,hidelinks,bookmarks=false]{hyperref}
\newcommand{\ie}{i.e.\ }
\newcommand{\cf}{cf.\ }
\newcommand{\resp}{respectively\ }
\newcommand{\Subsection}{\subsection}
\providecommand{\nobreakdash}{\nobreak\hbox{-}}
\setlength{\emergencystretch}{2em}
\multlinegap0pt
\usepackage{amssymb}
\usepackage[shortlabels]{enumitem}
\usepackage[normalem]{ulem}
\usepackage{xcolor}
\newtheorem{setup}{Setup}
\newtheorem{lemma}{Lemma}
\usepackage{cleveref}
\crefname{setup}{Setup}{Setups}
\crefname{lemma}{Lemma}{Lemmas}
\title{A blowup criterion for regularity}
\author{Pat Lank}
\date{Source: arXiv:2609.15691v1 (CC BY 4.0)}
\begin{document}
\maketitle

\noindent\emph{Source and licence.} A blowup criterion for regularity. Version arXiv:2609.15691v1 (CC BY 4.0), distributed by arXiv under the Creative Commons Attribution 4.0 International licence, http://creativecommons.org/licenses/by/4.0/, as recorded in the arXiv licence metadata for this exact version and shown on the abstract page https://arxiv.org/abs/2609.15691v1. Full licence notice: \texttt{licenses/arxiv-2609.15691-CC-BY-4.0.txt}.\par\smallskip

\noindent\textit{Editorial scope.} Counted unit: the article's lemma lem:ring\_presentationJona3 (source0.tex lines 638--647). The article's complete original proof of it is delivered with it (source0.tex lines 649--667, one \texttt{proof} environment of 19 lines). No mathematics is written, repaired or re-derived: the statement and the proof are the article's own lines, in the article's own notation, and no retained line is substituted or re-flowed.  Retained uncounted as the source-local context the proof needs: lines 299--307.  Nothing was omitted from any retained span, so the package carries no omission record.  No retained line contains a citation, so the wrapper carries no bibliography and no bibliography entry.  No retained line contains a figure environment, an image-inclusion command or drawing code, so no image is delivered and the package declares no asset dependency.  Editorial formatting only, in addition: an AMS \texttt{amsart} preamble that restates the article's own declarations and macros used by the retained lines, namely the environments \texttt{setup}, \texttt{lemma} on separate counters, together with the standard packages those lines need.  The retained environments print in this document as Setup 1, Lemma 1 in order of appearance, whereas the article prints the same items with the numbering of its own full document.  The complete native source of the article as distributed in the arXiv e-print for this exact version is bundled unchanged as \texttt{sources/210440/source0.tex}.
\begin{setup}
    \label{setup:main}
    Suppose $(R,\mathfrak{m},k)$ is a Noetherian local domain which admits an $R$-regular sequence $a,b\in \mathfrak{m}$. 
    %This means $\dim R \geq \operatorname{depth}(R) \geq 2$.
    Let $X:=\operatorname{Spec}(R)$, $i \colon \operatorname{Spec}(k) \to X$ be the associated closed immersion, and $c:=a^2 b$.
    Choose a minimal set of generators $r_1,\ldots,r_e$ for $\mathfrak{m}$.
    Define the ideals $P=(a,b)^3$ and $J=(a^3,ab^2,b^3)+ a^2 b \mathfrak{m}$ of $R$.
    Denote by $\mathcal{P}$ and $\mathcal{J}$ respectively for the sheafifications of $P$ and $J$.
\end{setup}

\begin{lemma}
    \label{lem:ring_presentationJona3}
    Assume \Cref{setup:main}. 
    Let $T$ be an indeterminate. 
    There exists a surjective ring homomorphism $\phi\colon R[T] \to R[\frac{b}{a}]$ given by $T\mapsto \frac{b}{a}$.
    Moreover, the kernel of $\phi$ is the ideal $(aT - b)$, and hence,
    \begin{displaymath}
        R[T]/(aT - b) \cong R[\frac{b}{a}].
    \end{displaymath}
\end{lemma}

\begin{proof}
    The mapping is a well-defined ring homomorphism and is surjective by construction.
    Since $(aT-b) \subseteq \ker (\phi)$, we need to check the reverse inclusion. 
    Fix $F(T)\in \ker (\phi)$. 
    If $F(T)=0$, then there is nothing to check.
    Assume $F(T)\not=0$.
    Set $n:= \operatorname{deg} (F(T))$.
    We prove the claim by induction on $n$. 
    If $n=0$, then $F(T)\in R$ is mapped to zero in $\operatorname{Frac}(R)$, which implies $F(T)=0$.
    Now assume that $n\geq 1$.
    Write $F(T)=\sum^n_{i=0} r_i T^i$ with $r_n\not=0$.
    As $\sum^n_{i=0} r_i (\frac{b}{a})^i = 0$, multiplying by $a^n$ yields $\sum^n_{i=0} a^{n-i} r_i b^i = 0 \in R$. 
    Reducing modulo the principal ideal $(a)$, we obtain that $r_n b^n \in (a)$.
    However, $a$ and $b$ form an $R$-regular sequence, and so $r_n \in (a)$. 
    Then $r_n = s_n a$ for some $s_n\in R$. 
    It follows that $F(T)- s_n T^{n-1}(a T - b) \in \ker (\phi)$ and has degree strictly less than $n$. 
    Thus, by the induction step, $F(T)- s_n T^{n-1}(a T - b)$ is in the ideal $(a T - b)$. 
    Since $s_n T^{n-1}(a T - b)$ is in the ideal $(a T - b)$, it follows that $F(T)\in (a T - b)$.
\end{proof}

\end{document}
